\begin{proof}[Proof of
\cref{prop:p082-formal-twistor-counterexample}]
Put
\[
                 R=\mathbf C[[t]],\qquad K=\mathbf C((t)).
\]
We separate the four assertions because the passage from formal to analytic
line bundles is load-bearing.

\emph{The analytic space.}
In the terminology of \cite[\S2.1]{HMS09Orientation}, a formal deformation
of $S$ is a smooth proper formal $R$-scheme whose special fibre is
$S$.  The formal twistor space associated to $\omega$ is such a
deformation.  Its Raynaud generic fibre is therefore a smooth proper strict
$K$-analytic space.  More precisely, the generic-fibre construction in
\cite[\S5.2.2]{Temkin12Introduction} produces a strict analytic space from
an admissible formal scheme of finite type, and
\cite[Fact~5.2.2.4]{Temkin12Introduction} takes formal smoothness to
rig-smoothness and formal properness to analytic properness.  A proper
analytic space is boundaryless, so rig-smoothness here is smoothness.  See
also \cite[Introduction and Remark~2.2]{HMS11FormalFibre}.  The original
formal model $\mathfrak X$ has projective special fibre $S$.

We verify connectedness rather than inferring it merely from domains on
formal affine charts.  Write
\[
  \mathfrak X_n
   =\mathfrak X\times_{\operatorname{Spf}R}\operatorname{Spec}
        \bigl(R/(t^{n+1})\bigr).
\]
Flatness gives, for $n\geq 1$, an exact sequence
\begin{equation}\label{eq:p082-infinitesimal-structure-sequence-candidate}
  0\longrightarrow \mathcal O_S
   \xrightarrow{\ t^n\ }
   \mathcal O_{\mathfrak X_n}
   \longrightarrow \mathcal O_{\mathfrak X_{n-1}}
   \longrightarrow 0.
\end{equation}
For a K3 surface,
$H^0(S,\mathcal O_S)=\mathbf C$ and
$H^1(S,\mathcal O_S)=0$.  Taking global sections in
\eqref{eq:p082-infinitesimal-structure-sequence-candidate} and comparing
with
\[
  0\longrightarrow \mathbf C\,t^n
   \longrightarrow R/(t^{n+1})
   \longrightarrow R/(t^n)
   \longrightarrow 0
\]
shows inductively that the structural map is an isomorphism
\begin{equation}\label{eq:p082-infinitesimal-global-functions-candidate}
       R/(t^{n+1})
          \xrightarrow{\ \sim\ }
       H^0(\mathfrak X_n,\mathcal O_{\mathfrak X_n})
       \qquad(n\geq0).
\end{equation}
Indeed, the right-hand restriction map is surjective because the next term
in the cohomology sequence is $H^1(S,\mathcal O_S)=0$, and the kernels in
the two displayed short exact sequences are identified by $t^n\mapsto
t^n\cdot1$.  Passing to inverse limits gives
\begin{equation}\label{eq:p082-formal-global-functions-candidate}
                \Gamma(\mathfrak X,\mathcal O_{\mathfrak X})=R.
\end{equation}
Here global sections commute with the inverse limit defining the structure
sheaf of the noetherian formal scheme.

By \cite[Remark~2.2 and Proposition~2.3]{HMS11FormalFibre}, coherent
sheaves on the Raynaud generic fibre are the categorical generic fibres of
coherent formal sheaves, and their morphisms are obtained by tensoring with
$K$.  Applying this to the structure sheaf and using
\eqref{eq:p082-formal-global-functions-candidate} yields
\begin{equation}\label{eq:p082-analytic-global-functions-candidate}
             \Gamma(X,\mathcal O_X)
              =\operatorname{End}(\mathcal O_X)
              \simeq R\otimes_R K=K.
\end{equation}
A disconnection of $X$ would give a nontrivial idempotent in
$\Gamma(X,\mathcal O_X)$, whereas the field $K$ has no such idempotent.
Thus $X$ is connected.  This proves the first assertion.

\emph{The analytic Picard group.}
Let $L$ be an analytic line bundle on $X$.  The fixed-model existence
theorem \cite[Theorem~B.1.8(a)]{CLD25} gives a coherent
$\mathcal O_{\mathfrak X}$-module $\mathcal F_0$ whose generic fibre is
$L$.  Let $\mathcal T$ be its $t$-power-torsion subsheaf and put
\[
                    \mathcal F=\mathcal F_0/\mathcal T.
\]
Noetherianity makes $\mathcal T$ coherent and locally annihilated by one
power of $t$.  Hence $\mathcal F_\eta=L$, while multiplication by $t$
is injective on $\mathcal F$.

We claim that $\mathcal F$ is torsion-free of rank one as an
$\mathcal O_{\mathfrak X}$-module.  Work in a formal local ring $A$ of
$\mathfrak X$.  The local map $R\to A$ is flat because
$\mathfrak X\to\operatorname{Spf}R$ is smooth, and $A/tA$ is a regular
local ring because it is a local ring of the smooth surface $S$.
Since $R$ is regular, \cite[Tag~031E]{StacksProject} shows that $A$ is
regular.  In particular, it is a local domain.  If
$0\neq a\in A$, $m\in\mathcal F_A$, and $am=0$, then after inverting
$t$ the element $a/1$ remains nonzero and the generic-fibre module is
locally free of rank one.  Therefore $m/1=0$, so $t^N m=0$ for some
$N$.  The $t$-torsion-freeness of $\mathcal F_A$ gives $m=0$.
Its rank is one because its generic fibre is $L$.

Consider the reflexive hull
\[
             \mathcal M:=\mathcal F^{\vee\vee}.
\]
It is coherent and has rank one because finite duals preserve coherence and
localization at the fraction field identifies its generic fibre with a
one-dimensional vector space.  It is reflexive by
\cite[Tag~0AV3]{StacksProject}, applied locally to
$\operatorname{Hom}_A(\mathcal F_A^\vee,A)$, and it can equivalently be
recognized by the normal-domain criteria in
\cite[Tags~0AVB and 0AVC]{StacksProject}.  To spell out the local algebra,
over a normal formal local domain $A$ the torsion-free rank-one module
$\mathcal F_A$ embeds in its one-dimensional vector-space generic fibre
and is therefore a fractional ideal.  Its double dual is the associated
divisorial fractional ideal.  Such
rank-one reflexive modules
are classified by the local Weil divisor class group
\cite[Tag~0EBM]{StacksProject}.  Every formal local ring $A$ of
$\mathfrak X$ is regular and hence a unique factorization domain
\cite[Tag~0AG0]{StacksProject}; its divisor class group is zero.
Consequently every stalk of $\mathcal M$ is free of rank one.  Coherence
then shows that $\mathcal M$ is a line bundle on $\mathfrak X$.

Formation of the dual of a finite module commutes with localization.
Since $\mathcal F[1/t]$ is already locally free, the natural map
$\mathcal F\to\mathcal M$ becomes an isomorphism after inverting $t$.
Consequently
\begin{equation}\label{eq:p082-double-dual-generic-fibre-candidate}
                         \mathcal M_\eta\simeq L.
\end{equation}
For a very general formal twistor deformation, every line bundle on
$\mathfrak X$ is trivial by
\cite[Proposition~2.10(i)]{HMS09Orientation}.  Thus
$\mathcal M\simeq\mathcal O_{\mathfrak X}$, and
\eqref{eq:p082-double-dual-generic-fibre-candidate} gives
$L\simeq\mathcal O_X$.  Since $L$ was arbitrary, this proves
\eqref{eq:p082-formal-twistor-picard-zero-candidate}.

\emph{A nonconstant formal metric on the trivial bundle.}
Choose a closed point $p\in S$, and let
\[
 \mathcal J=\ker\bigl(
       \mathcal O_{\mathfrak X}\longrightarrow (i_p)_*\mathbf C
                         \bigr).
\]
This is a coherent open ideal because it contains $t$.  Let
\[
       \mu\colon\mathfrak X'
          =\operatorname{Bl}_{\mathcal J}(\mathfrak X)
          \longrightarrow\mathfrak X
\]
be the admissible formal blow-up.  Since
$\mathcal J[1/t]=\mathcal O_{\mathfrak X}[1/t]$, its generic-fibre map is
an isomorphism by Raynaud's localization theorem
\cite[Exercise~5.3.2.3 and Fact~5.3.3.1]{Temkin12Introduction}; see also
the direct chart proof
in \cite[\S8.4, Proposition~2]{Bosch14FormalRigid}.  The special-fibre morphism
$\mathfrak X'_s\to S$ is a relative $\operatorname{Proj}$ of the
reduction of the finitely generated Rees algebra of $\mathcal J$
\cite[Tags~052P, 0804, and 01OF]{StacksProject};
hence it is projective.  Since $S$ is projective,
$\mathfrak X'_s$ is projective as well.  Thus even the formal model that
defines the metric below has projective special fibre.

The tautological transform of the centre ideal is invertible; write
\begin{equation}\label{eq:p082-exceptional-line-bundle-candidate}
       \mathcal J\mathcal O_{\mathfrak X'}
                 =\mathcal O_{\mathfrak X'}(-E).
\end{equation}
The inclusion of this ideal in $\mathcal O_{\mathfrak X'}$ identifies the
generic fibre of \eqref{eq:p082-exceptional-line-bundle-candidate} with
$\mathcal O_X$.  Denote its canonical formal metric by
$\lVert\cdot\rVert_E$.  By
\cite[Proposition~17.7.5(b)]{CLD25}, it is a strictly integral parlinear
metric.

We compute both the special fibre and the metric in one chart.  Smoothness
of relative dimension two gives a noncanonical identification
\[
       \widehat{\mathcal O}_{\mathfrak X,p}
             \simeq \mathbf C[[t,x,y]],
       \qquad
       \mathcal J_p=(t,x,y).
\]
Here the completion is the maximal-ideal completion of the noetherian local
ring $\mathcal O_{\mathfrak X,p}$, and the completion map is faithfully
flat by \cite[Tag~00MC]{StacksProject}.  Scheme-theoretic blow-up commutes
with this flat base change by \cite[Lemma~31.33.3, Tag~01OF]{StacksProject}.
Moreover, the formal blow-up is the $(t)$-adic completion of the
scheme-theoretic blow-up by
\cite[\S8.2, Proposition~6]{Bosch14FormalRigid}.  Consequently the affine
blow-up charts and their special-fibre equations may be computed after
passing to $\mathbf C[[t,x,y]]$.  Since the passage is faithfully flat,
the reduced component equations and the multiplicity-one assertions below
are detected there; adic completion does not alter the reduction modulo
$t$.
On the $x$-chart of the blow-up put
\[
                      T=t/x,\qquad Y=y/x.
\]
Thus $t=xT$, $y=xY$, and
$\mathcal J\mathcal O_{\mathfrak X'}=(x)$ on this chart.  Its special
fibre has equation
\[
                         xT=0.
\]
The exceptional component is $E=(x)$, while the strict transform of
$S$ is $S'=(T)$.  Both occur with multiplicity one in
\begin{equation}\label{eq:p082-special-fibre-components-candidate}
                         \operatorname{div}(t)=S'+E.
\end{equation}
Globally, $\mathfrak X'_s$ is the union of
$\operatorname{Bl}_p S$ and $E\simeq\mathbf P^2_{\mathbf C}$, meeting
along $\mathbf P^1_{\mathbf C}$.  Let $x_{S'}$ and $x_E$ be the
divisorial points of $X$ associated to these two components.

At the generic point of $S'$, the centre ideal $\mathcal J$ is the unit
ideal, since that generic point maps outside the closed point $p$.
Hence the ideal in
\eqref{eq:p082-exceptional-line-bundle-candidate} is generated by $1$,
and the formal-frame convention gives
\begin{equation}\label{eq:p082-metric-value-strict-transform-candidate}
                         \lVert1(x_{S'})\rVert_E=1.
\end{equation}
At the generic point of $E$, the function $x$ is a formal frame of
$\mathcal O_{\mathfrak X'}(-E)$.  Under the chosen generic-fibre
identification the trivial section satisfies
\[
                              1=x^{-1}\,x.
\]
Therefore
\begin{equation}\label{eq:p082-metric-value-exceptional-candidate}
                         \lVert1(x_E)\rVert_E=|x(x_E)|^{-1}.
\end{equation}
For a component $V$ of a formal special fibre, its divisorial seminorm is
normalized by
\[
 |f(x_V)|=|t|^{
       \operatorname{ord}_V(f)/\operatorname{ord}_V(t)}.
\]
Since $t=xT$ and $T$ is a unit at the generic point of $E$,
$\operatorname{ord}_E(x)=\operatorname{ord}_E(t)=1$.  It follows that
\begin{equation}\label{eq:p082-two-metric-values-candidate}
        \lVert1(x_E)\rVert_E=|t|^{-1}\neq1
              =\lVert1(x_{S'})\rVert_E.
\end{equation}
For complete sign bookkeeping, the logarithmic norm
$\log\lVert1\rVert_E$ has respective values
$-\log|t|$ and $0$ at $x_E$ and $x_{S'}$.  With the CLD weight
convention $\varphi=\log\lVert1\rVert_E^{-1}$, the respective values are
$\log|t|$ and $0$.  This proves the third assertion.

\emph{Constancy of globally dpsh metrics on $\mathcal O_X$.}
We first show that every globally psh metric on an analytically trivial
line bundle over $X$ is constant in a global trivialization.  Fix a
nowhere-vanishing analytic section $s$.  By
\cite[\S17.8.5]{CLD25}, the metric is a uniform limit of smooth psh metrics
on the same line bundle.  If
\[
                         u_n=\log\lVert s\rVert_n^{-1}
\]
is the weight of such an approximating metric, then $u_n$ is a smooth
CLD-psh function.  CLD-psh functions are classically psh
\cite[\S1.2]{GR25Classical}.  The proper space $X$ is compact and
boundaryless, so $u_n$ attains a maximum.  Its maximum locus is closed by
continuity.  At each point of that locus, the local maximum principle
\cite[Theorem~6.1]{GR25Classical} makes $u_n$ constant on a
neighbourhood; hence the maximum locus is also open.  Connectedness of
$X$ makes it all of $X$.  Thus every $u_n$ is constant, and uniform
convergence shows that the limiting weight is constant as well.

Let now $\lVert\cdot\rVert$ be a globally dpsh metric on
$\mathcal O_X$.  By the definition in \cite[\S17.8.5]{CLD25}, there are
globally psh metrized line bundles
$(L_1,\lVert\cdot\rVert_1)$ and
$(L_2,\lVert\cdot\rVert_2)$, together with an isometric isomorphism
\begin{equation}\label{eq:p082-global-dpsh-decomposition-candidate}
  (\mathcal O_X,\lVert\cdot\rVert)
       \simeq
  (L_1,\lVert\cdot\rVert_1)
       \otimes (L_2,\lVert\cdot\rVert_2)^{-1}.
\end{equation}
By \eqref{eq:p082-formal-twistor-picard-zero-candidate}, choose analytic
trivializations $s_i$ of $L_i$.  The preceding paragraph shows that
\[
                        u_i=\log\lVert s_i\rVert_i^{-1}
\]
is constant for $i=1,2$.  In these trivializations, the isomorphism in
\eqref{eq:p082-global-dpsh-decomposition-candidate} sends $1$ to
\[
                         f\,s_1\otimes s_2^{-1}
\]
for a global analytic unit
$f\in\Gamma(X,\mathcal O_X^\times)$.  Hence the weight
$u=\log\lVert1\rVert^{-1}$ is
\begin{equation}\label{eq:p082-global-dpsh-weight-candidate}
                         u=u_1-u_2-\log|f|.
\end{equation}
The function $\log|f|$ is classically pluriharmonic by
\cite[Proposition~7.2(2)]{GR25Classical}.  It attains a maximum on the
compact space $X$.  Its maximum locus is open by the pluriharmonic local
maximum principle, closed because $\log|f|$ is continuous, and nonempty;
connectedness makes
it all of $X$.  Thus $\log|f|$ is constant.  Equation
\eqref{eq:p082-global-dpsh-weight-candidate} proves that $u$ is constant,
which establishes the fourth assertion.

The formal metric $\lVert\cdot\rVert_E$ is nonconstant by
\eqref{eq:p082-two-metric-values-candidate}, whereas every globally dpsh
metric on $\mathcal O_X$ is constant.  It is therefore not globally
dpsh, which disproves the proposed extension of
\cite[Corollary~17.8.9]{CLD25}.

This conclusion concerns only the second paragraph of
\cite[Remark~17.8.10]{CLD25}.  The constructed metric is parlinear; no
claim that it is psh is made.  In particular, the argument does not answer
the first paragraph's continuous-psh-to-global-psh question, nor does it
identify global dpsh approximation with local psh approximation.
\end{proof}
